% =============================================================================
% Copyright (c) 2026 Mert Torun, M.Sc. (mtorun0x7cd) <info@mtorun0x7cd.com>
% 
% Permission is hereby granted, free of charge, to any person obtaining a copy
% of this software and associated documentation files (the "Software"), to deal
% in the Software without restriction, including without limitation the rights
% to use, copy, modify, merge, publish, distribute, sublicense, and/or sell
% copies of the Software, and to permit persons to whom the Software is
% furnished to do so, subject to the following conditions:
% 
% The above copyright notice and this permission notice shall be included in
% all copies or substantial portions of the Software.
% 
% THE SOFTWARE IS PROVIDED "AS IS", WITHOUT WARRANTY OF ANY KIND, EXPRESS OR
% IMPLIED, INCLUDING BUT NOT LIMITED TO THE WARRANTIES OF MERCHANTABILITY,
% FITNESS FOR A PARTICULAR PURPOSE AND NONINFRINGEMENT. IN NO EVENT SHALL THE
% AUTHORS OR COPYRIGHT HOLDERS BE LIABLE FOR ANY CLAIM, DAMAGES OR OTHER
% LIABILITY, WHETHER IN AN ACTION OF CONTRACT, TORT OR OTHERWISE, ARISING FROM,
% OUT OF OR IN CONNECTION WITH THE SOFTWARE OR THE USE OR OTHER DEALINGS IN THE
% SOFTWARE.
% =============================================================================

\chapter{Conclusion}
\label{ch:conclusion}
This chapter summarizes the findings of this work and outlines final conclusions and future work.

\section{Summary}
This thesis presented a modular integration architecture for containerized FPGA toolchains.
\begin{itemize}
    \item \textbf{Contributions:} A modular Hybrid Strategy Pattern for transparent, opt-in container dispatch; cross-platform permission and transport handling (host UID/GID injection, Unix-socket and named-pipe transports); and a toolchain-image pinning mechanism for reproducible build environments.
    \item \textbf{Result:} The solution eliminates ``Environment Drift'' while maintaining native IDE responsiveness. The overhead the extension adds beyond raw containerization is small and dominated by a fixed per-invocation startup cost that amortizes as design size grows; under x86-on-ARM64 emulation an additional translation tax applies. Native \texttt{x86\_64} hosts (Machine~3, Windows, and Machine~2, Linux) reproduced the same fixed-floor structure free of emulation, and the same-build native baseline on Machine~2 fixed the residual container-versus-native factor at $\alpha \approx 1$---a fixed lifecycle floor with no per-instruction tax---completing the overhead picture across all three hosts. Crucially, across the seven-workload FPGA corpus---iCE40 and ECP5 synthesis, place-and-route and bitstream packing, plus Icarus~Verilog simulation---every workload executed to success through the hermetic container, where a native toolchain would be subject to package drift, so the container makes toolchain execution independent of the host's package environment and resolves cascading dependency conflicts.
\end{itemize}

\section{Evaluation Summary and Key Architectural Insights}
The central empirical insight of this work is a decomposition of the containerized execution cost rather than a single overhead figure. The cost of an invocation is the sum of a container lifecycle floor---container creation, namespace setup, volume mounting, and stream attachment, measured at roughly \SI{0.25}{\second} for the raw CLI backend on Machine~1 (\SI{0.33}{\second} and \SI{0.42}{\second} on Machines~2 and~3)---and a per-instruction execution tax that is unity for native-architecture execution and greater than unity under binary translation. On top of that floor, the extension's strategy dispatch adds only a bounded, fixed marginal cost---roughly \SI{0.18}{\second} per invocation on Machine~1, tightly clustered across all seven FPGA workloads and matching the bare no-op floor delta---and no measurable marginal compute overhead, so its relative cost is large only for sub-second tasks and amortizes toward negligibility as compilation time grows. On Apple Silicon the toolchain image executes under Rosetta~2, so any native-versus-container comparison on Machine~1 would conflate containerization with instruction-set emulation; the translation tax itself is small and approximately constant at scale. Native \texttt{x86\_64} measurements on Machine~3 (Windows) and Machine~2 (Linux) confirmed the fixed-floor structure without emulation, and the same-build native baseline on Machine~2---the one host where the native and containerized toolchains are the identical build---fixed the container-versus-native factor at $\alpha \approx 1$: the additive container floor stays near-constant at \SIrange{0.26}{0.33}{\second} whether the native work beneath it is \SI{16}{\milli\second} or half a second, so containerization adds a fixed lifecycle cost and no per-instruction tax. A systematic comparison of container runtimes (Docker Desktop, Podman, OrbStack, and the WSL2 backend) is left to future work, as the present measurements were collected on a single runtime per host.

As a cross-platform robustness detail, the Windows named-pipe transport requests \texttt{TokenImpersonationLevel.Identification} rather than the default impersonation level, so a daemon endpoint can verify the connecting client's identity without being able to act on its behalf. This, together with the remaining container-hardening defaults---non-root execution, dropped Linux capabilities, \texttt{no-new-privileges}, a process-count cap, and shell-less argument execution---is a design property of the architecture. A dedicated threat-model analysis of these mechanisms is the subject of separate work and is deliberately out of scope for this thesis.

To ensure observability without data corruption, the telemetry pipeline separates two write paths. High-frequency error events are enqueued lock-free into an unbounded channel and drained by a single background consumer, decoupling the calling thread from disk I/O. The execution-log file write, by contrast, is serialized across threads and processes by a reader-writer lock (\path{ReaderWriterLockSlim}) and a system-wide mutex (\path{System.Threading.Mutex}), so that concurrent compilation steps cannot interleave or corrupt the structured log. This division keeps the hot path non-blocking while guaranteeing that persisted telemetry is cleanly serialized.

\section{Future Work}
Future development on the OneWare Container Extension will focus on two primary architectural enhancements:
\begin{itemize}
    \item \textbf{Multi-Architecture OCI Registry Manifests:} Optimizing cross-architecture performance on ARM64 host machines by dynamically querying OCI Registry Manifests to verify the existence of native \texttt{linux/arm64} image variants before falling back to Rosetta 2 translation or QEMU emulation.
    \item \textbf{Hardware-in-the-Loop (HIL) Sandboxing:} Integrating USB device forwarding mechanisms (e.g., forwarding FTDI or UART devices via \texttt{usbipd} or native Linux device nodes) to allow programmer tools like \texttt{openFPGALoader} and \texttt{iceprog} to deploy compiled bitstreams directly onto target FPGA boards securely from within the sandbox.
\end{itemize}
